\section{The native original order is the actual stalk quotient length}

\begin{theorem}[Actual local order and true quotient length]
\label{thm:fc-scheme-order-local-length}
\leanok
\lean{EnvelopeProperty.SchemeOrderLocalLength.principal_stalk_quotient_finiteLength, EnvelopeProperty.SchemeOrderLocalLength.ord_germToFunctionField_eq_stalk_length}
On an actual integral locally Noetherian scheme, let $x$ have actual
codimension one and let $c$ be a nonzero section on an actual open
neighborhood of $x$. Its actual stalk quotient by the principal ideal
of its original germ has finite length. The native original order
$\operatorname{ord}_x(c)$ in the actual original function field equals
this actual length, regarded as an integer.
\end{theorem}
\begin{proof}
The original germ map is injective by integrality, so the original germ
is a nonzerodivisor in the actual original stalk domain. Its stalk
has dimension at most one by the actual codimension-one hypothesis.
The true Noetherian principal quotient length theorem therefore proves
finite length. The original native function-field germ is the image
of the same stalk germ under the actual fraction-field map. The native
fraction-field order extends the actual stalk order, which is defined
as precisely this actual principal quotient length. Finiteness makes
its natural and integer values agree, giving the equality.
\end{proof}

\begin{theorem}[Actual local rational orders are differences of true lengths]
\label{thm:fc-scheme-rational-order-local-length}
\leanok
\lean{EnvelopeProperty.SchemeOrderLocalLength.ord_div_germToFunctionField_eq_stalk_lengths}
\uses{thm:fc-scheme-order-local-length}
For two such nonzero actual sections $c,d$ on the same original open,
the native order of their actual original function-field quotient is
the difference of their actual principal stalk quotient lengths.
\end{theorem}
\begin{proof}
The actual original function-field germ is injective, so both elements
and their quotient are nonzero. Apply native order multiplicativity
to $(c/d)d=c$, cancel in the integers, and substitute the preceding
same-stalk length identities for $c$ and $d$.
\end{proof}
